\documentclass[12pt,a4paper,oneside,naustrian]{amsbook}
\usepackage[T1]{fontenc}
\usepackage[utf8]{inputenc}
\setcounter{secnumdepth}{2}
\setcounter{tocdepth}{1}
\usepackage{parskip}
\usepackage{amsthm}
\usepackage{amssymb}
\usepackage{graphicx}
\usepackage{setspace}
\usepackage{natbib}
\usepackage{hyperref}

\sloppy

\onehalfspacing

\makeatletter

\pdfpageheight\paperheight
\pdfpagewidth\paperwidth


%%%%%%%%%%%%%%%%%%%%%%%%%%%%%% Textclass specific LaTeX commands.
\numberwithin{section}{chapter}
\numberwithin{equation}{section}
\numberwithin{figure}{section}
\theoremstyle{plain}
\newtheorem{thm}{\protect\theoremname}
\theoremstyle{plain}
\newtheorem{prop}[thm]{\protect\propositionname}

%%%%%%%%%%%%%%%%%%%%%%%%%%%%%% User specified LaTeX commands.
\usepackage{graphicx,lipsum}
\usepackage[total={20.9cm,29.7cm},
margin=2.5cm,top=1.5cm,includeheadfoot]{geometry}
\setlength{\footskip}{2\baselineskip}

\makeatother

\usepackage{babel}
\providecommand{\propositionname}{Satz}
\providecommand{\theoremname}{Theorem}

\begin{document}

\savegeometry{saved}
\leftmargin 0cm
\textwidth 16.7cm
\textheight 23cm
\topmargin -2.7cm
\oddsidemargin 0.1cm
\parindent 0pt

\def\title{4D Druck im Unterricht:\\ Fraktale in der Unterstufe}
\def\name{A.~B.~Conrad Dachsteiner}
\def\institute{Institut für \\Didaktik der Mathematik}
\def\supervisor{Dr.~Zoltán Kovács}
\def\assist{Dr.~Zsolt Lavicza}
\def\date{Juni 2025}

\thispagestyle{empty}
\def\ifundefined#1{\expandafter\ifx\csname#1\endcsname\relax}
\DeclareFontShape{OT1}{cmss}{m}{n}
  {<5><6><7><8><9><10><10.95><12><14.4><17.28><20.74><24.88><29.86><35.83><42.99><51.59><67><77.38>cmss10}{}
\DeclareFontShape{OT1}{cmss}{bx}{n}
  {<5><6><7><8><9><10><10.95><12><14.4><17.28><20.74><24.88><29.86><35.83><42.99><51.59><67><77.38>cmssbx10}{}
\makeatletter
\def\Huge{\@setfontsize\Huge{29.86pt}{36}}
\makeatother
\unitlength 1cm
\sffamily
\begin{picture}(16.7,0)\put(-1.0,-2.5){

\includegraphics[width=7.22cm]{PHDL}\hspace{6mm}\includegraphics[width=3.41cm]{PH-OOe}\hspace{13mm}\includegraphics[width=4.94cm]{JKU}

}
\put(12.5,-4.2){\begin{minipage}[t]{5cm}\footnotesize%
 Eingereicht von\\ {\bfseries\name}%
\vskip 4mm%
 Angefertigt am\\ {\bfseries\institute}%
\vskip 4mm%
 Beurteiler / Beurteilerin\\ {\bfseries\supervisor}%
\vskip 4mm%
 Mitbetreuung\\ {\bfseries\assist}%
\vskip 4mm%
\date
\end{minipage}}

\put(0,-25){\begin{minipage}[t]{3.9cm}\footnotesize%
{\bfseries Private P\"adagogische \\ Hochschule der \\ Di\"ozese Linz}\\
Salesianumweg 3\\ 4020 Linz, Österreich\\ phdl.at
\end{minipage}}

\put(6.5,-25){\begin{minipage}[t]{5cm}\footnotesize%
{\bfseries Pädagogische Hochschule\\ Oberösterreich}\\
Kaplanhofstraße 40 /\\ Huemerstrasse 3 - 5\\
4020 Linz, Österreich\\ ph-ooe.at
\end{minipage}}

\put(12.9,-25){\begin{minipage}[t]{3.9cm}\footnotesize%
{\bfseries JOHANNES KEPLER\\ UNIVERSITÄT LINZ}\\
Altenbergerstra{\ss}e 69\\ 4040 Linz, Österreich\\
www.jku.at\\ DVR 0093696
\end{minipage}}

\put(0,-15.2){\begin{minipage}[b]{15cm}\Huge\bfseries\title\end{minipage}}

\put(0,-18.3){\begin{minipage}[t]{12cm}%
 {\large Bachelorarbeit}%
 \vskip 2mm%
 zur Erlangung des akademischen Grades%
 \vskip 3mm%
 {\large Bachelor of Education}
 \vskip 3mm%
 im Bachelorstudium
\vskip 3mm%
{\large Lehramt Sekundarstufe (Allgemeinbildung)}
\end{minipage}}
\end{picture}
\eject
\rmfamily

\loadgeometry{saved}
\pagestyle{empty}
\newpage{}
\begin{center}
\vspace*{0.3\paperheight}
Für meine geliebte Mama
\par\end{center}

\chapter*{Kurzfassung}

\lipsum[1]

\vskip2cm

\begingroup
\renewcommand{\cleardoublepage}{}
\renewcommand{\clearpage}{}
\chapter*{Abstract}
\lipsum[2]
\endgroup

\begin{center}
\thispagestyle{empty}\newpage\tableofcontents{}
\par\end{center}

\chapter{Einführung}

Bla\ldots{} bla\ldots{} bla\ldots{} \lipsum[1-10]

\section{Wichtige Bemerkungen}

\lipsum[11-20]

\section{Weitere wichtige Bemerkungen}

\lipsum[21-30]

\chapter{Tiefes Wasser}

In diesem Kapitel\ldots{} blabla\ldots{} bla-bla\ldots{} \lipsum[31]

\section{Sehr wichtige Bemerkungen}

\lipsum[33-35]

Nun schließen wir darauf, dass
\[
\sum_{n=0}^{\infty}\frac{1}{n^{2}}=\frac{\pi^{2}}{6}.
\]
Unsere Bemerkungen können in dieser Form zusammengefasst:
\begin{prop}
Für jedes $k\in\mathbb{N}$, $k>1$ gilt

\[
\sum_{n=0}^{\infty}\frac{1}{n^{k}}<\infty.
\]

\begin{proof}
Trivial.
\end{proof}
\end{prop}

Ein weiterer Beweis befindet sich im \citet{Adams1980}.
\begin{thebibliography}{}
\bibitem[Adams(1980)]{Adams1980}Douglas Adams: \emph{Das Restaurant am Ende
des Universums} \emph{(The Restaurant at the End of the Universe)},
ISBN 3-453-14698-0. Originalveröffentlichung 1980. Veröffentlichung
der deutschen Übersetzung 1983.

\end{thebibliography}

\appendix % Anhang

% ---------------------------------------------------------
\chapter*{Erklärung} % Eidesstattliche Erklärung
% ---------------------------------------------------------

Hiermit erkläre ich an Eides statt, dass ich diese Arbeit eigenständig
verfasst habe. Direkte und indirekte Zitate (Texte als auch Abbildungen)
sind unter Angabe der Quelle nach den Regeln des wissenschaftlichen
Zitierens und der guten wissenschaftlichen Praxis kenntlich gemacht.
Dies schließt auch verwendete Abbildungen ein.

Bei Verwendung von Künstlicher Intelligenz sind die entsprechenden
Anwendungen mit ihrem Produktnamen und der in dieser Arbeit verwendeten
Funktionen im Folgenden vollständig dokumentiert.

Ich bin mir darüber im Klaren, dass ich als Autor allein für den Inhalt
verantwortlich bin. Das Urheberrecht sowie alle damit verbundenen Rechte
an Text, Struktur, Inhalt und Bilder verbleiben bei mir.

Zur (sprachlichen) Überarbeitung wurden folgende Software- und KI-gestützte Anwendungen genutzt:

\begin{itemize}
    \item Overleaf: zur Überprüfung der Rechtschreibung
    \item Microsoft Word: zur Überprüfung der Rechtschreibung und Grammatik
    \item DeepL Write: zur stilistischen Optimierung von Formulierungen und Überprüfung der Rechtschreibung und Grammatik
    \item OpenAI ChatGPT: zur stilistischen Optimierung von Formulierungen und Überprüfung der Rechtschreibung und Grammatik
\end{itemize}


\end{document}
